\chapter{Discussion} \label{chap:discussion}  The ML and RK distributions, are displayed in Figure \ref{fig:theory_model} demonstrating the overall effect of geographic uncertainty applied to each complex trend and integration into the selected XGRB-RK method. \begin{figure}[h!]
    \centering
        \centering
        \includegraphics[width=0.75\linewidth]{ThesisFilePkg//figs/theory_model_preds.png}
    \caption{Distribution of Literature-Based Features with XGRB-RK Predictions.} 
    \label{fig:theory_model}
\end{figure} This chapter is dedicated to a brief application of the best model discussed in Chapter \ref{rlts:imgsandouts}. To do this, the American Community Survey 5-Year (ACS5Y) data set was collected for 2019 to 2023 provided by the \textcite{united_states_census_bureau_arizona_2025}. Yearly income data and population counts are utilized to illustrate potential growth over time. Given the extensive nature of this project, depicting roughly 2,125 days would be disingenuous to Mother Nature should a printed copy ever be required. This section overlays five randomly selected months with census tract statistics to limit the total number to around 150. These statistics are also presented alongside first-of-month depictions created in ArcGIS Pro (ArcPRO), which provides highly detailed maps illustrating the rapid integration of DAMO$_3$ with modern GIScience computer programs. 

These improvements are illustrated using a basic county-based separation of monitors at the census tract resolution. Certain DAMO$_3$ values are displayed in sets of four within this chapter; the remaining values are available in full size in the Appendix following the last appearance of DAMO$_3$ features used in prediction after Map \ref{appmap:smarkfeats_2023}. The DAMO$_3$ maps that do not include census-based shading colors and are shown without monitoring locations to provide a clear representation of the predicted result with XGRB-RK. In addition, to conserve space, DAMO$_3$ only maps will not be shown in the appendix. Only feature maps and associated tri-variate colored maps will be included in the appendix, as these have only been partially displayed and represent $\frac{1}{15}$th the count of DAMO$_3$ only maps, or 93.3 percent fewer images.

It was essential to capture all seasons in this sample; random months for each season were selected, including two months from the summer season. Once a month was selected, a year was randomly assigned to it. January 2019, October 2020, and April 2023 represent winter, fall, and spring, respectively, while July 2021 and June 2022 represent summer for the AOI. The colors used for the maps, distributions of DAMO$_3$ across PHOTUC, and the enhancements made by RK are discussed, concluding with future directions and general comments in Chapter \ref{chap:clcsn}. This SMaRK combination was chosen due to its complexity and effectiveness. The literature-based dataset (L.B.D) was utilized, as it was found to offer a broader variation of features compared to other datasets. Improvements of the RK methods in the AOI are also effectively demonstrated through MASO$_3$ values over time, split by county. 
%%%%%%%%%%%%%%%%%%%%%%%%%%%%%%%%%%%%%%%%%%%%%%%%%%%%%%%%%%%%%%%%
\section{Demographics in PHOTUC} \label{discussion:colordesc} As the demographic maps of the counties aim to demonstrate the functionality of Python and the implementation of raster data for the larger dataset, they will be reproduced at reduced sizes and referenced as Map \ref{appmap:azdem_2019}, \ref{appmap:azdem_2020}, A.n., etc. Combining the area of interest (AOI) with distributions of household income highlights various trends discussed in Chapter \ref{subs:sofat1}. Utilizing the theoretical dataset T.D alongside XGRB-RK, models predicting DAMO $_3$ concentrations are shown with the gathered basic population and income metrics. This illustrates high-resolution surface O$_3$ as it varies across a densely populated region. 

The maps of population distributions, income, and DAMO$_3$ for the AOI can be viewed further in Appendix \ref{appendixA}. The subsequent appendices are devoted to DAMO $_3$ predictions for the specified months and years mentioned earlier. A detailed breakdown of the results by county statistics offers visually compelling insights into the populations affected by DAMO$_3$ and the estimated extent of the impact, as determined by the SMaRK methodology. An example of the demographic coloring scheme is illustrated in Figure \ref{fig:azdem_2019} with the demographic statistics for the other four years available in the larger versions found in Appendix \ref{appmap:azdem_2020} to \ref{appmap:azdem_2023}. The boundary layers of census tract data are included in all maps and adhere to the same coloring scheme. 

The maps present an unconventional color scheme not typically seen in traditional GIScience. Counties are depicted in the same colors as shown by the solid black lines in Figure \ref{fig:theory_model}, while dashed lines indicate the boundaries. These statistical maps created in ArcGIS Pro utilize a bivariate color scheme comprised of yellow and blue. In this thesis, population counts were assigned as the total unweighted count within a census tract, represented by varying shades of yellow. Income is calculated as the weighted difference between mean and median income, mathematically represented as:
\[Display_{income}=\frac{(mean - median)}{mean}\]
where the standardized metric represents the skewness of total income for a census tract. Higher values denote a larger overall income and/or greater deviation from the mean. Some tracts with a high skew may be combined with tracts that generally exhibit higher values. This approach suggests that these counties would theoretically share similar growth potential, given their high base incomes. In summary, census tracts with darker shades of blue have a higher potential either split or high incomes with lighter shades representing equal to low average distributions.\begin{figure}[h!]
    \centering
    \includegraphics[width=0.8\linewidth]{ThesisFilePkg//figs/maps/azdem_2019.png}
    \caption{Bivariate Map of Census Tract Statistics for 2019}
    \label{fig:azdem_2019}
\end{figure} The equal combination of yellow and blue shades produces a green hue, signifying areas that maintain a balance between population density and income. This illustrates counties with significant growth potential. The textures for the bivariate maps are initially presented in grayscale. They are subsequently overlaid with DAMO$_3$ averages for each county derived from XGRB-RK in a continuous red color scheme of DAMO$_3$ values depicted later in \ref{fig:dem_o3_2019}. DAMO$_3$, compared with monitoring statistics, is represented in blue.  
%%%%%%%%%%%%%%%%%%%%%%%%%%%%%%%%%%%%%%%%%%%%%%%%%%%%%%%%%%%%%%%%%%%%%%%%%%%%%%%%%%%%%%%%%%%%%%%%%%%%%%%%%%%%%%%%%%%%
\section{DAMO$_3$ in Southern Arizona} \label{discussion:damo3soaz} While the core statistics underlying the coloring scheme may appear complex, they provide valuable insights by illustrating interesting spatial trends within PHOTUC. Figure \ref{fig:ozone_2023-04-06to09} illustrates a piece of the data set created for April, with DAMO$_3$ initially shaded blue. \begin{figure}[h!]
    \centering
    \includegraphics[width=0.35\linewidth]{ThesisFilePkg/figs/maps/ozone_2023-04-06.png}
    \includegraphics[width=0.35\linewidth]{ThesisFilePkg/figs/maps/ozone_2023-04-07.png}
    \includegraphics[width=0.35\linewidth]{ThesisFilePkg/figs/maps/ozone_2023-04-08.png}
    \includegraphics[width=0.35\linewidth]{ThesisFilePkg/figs/maps/ozone_2023-04-09.png}
    \caption{DAMO$_3$ in PHOTUC from April 6th, 2023 to April 9th, 2023}
    \label{fig:ozone_2023-04-06to09}
\end{figure} The numerous maps created for this section aim to briefly expand on the who, when, and where of DAMO$_3$ and the potential outcomes in areas exposed to high concentrations. Due to the low terrain variance, particularly around mountainous regions, many urban populations reside in basins that facilitate the accumulation of DAMO$_3$, which originates from higher elevations surrounding Phoenix and Tucson. O$_3$, as an oxidizing agent, has been found to participate in redox reactions while simultaneously reducing due to its molecular instability (e.g. \cite{marmett_role_2023}; \cite{qiu_enabling_2025}; \cite{yang_enhancing_2024}). Southeastern Arizona is characterized by a semi-arid climate, with about 300 mm of rain per year as per \textcite{le_tourneau_socio-ecological_2024}. 

The AOI experiences cycles of 3-month distributions with less than 5 mm of precipitation. These can be seen in \ref{fig:scatters} in Chapter \ref{tmas:findatasets}, where precipitation was essentially unusable. The average annual temperature exceeds 294 K (~20\textdegree C), a common characteristic in the western US at 30\textdegree degrees latitude (\cite{marshall2010}). Desert-like climates in between high-elevation areas typical of the "basin and range" type landscape, allow for high biodiversity in "hotspots" throughout PHOTUC. Mountain ranges, with peaks up to 2800 m, are called "sky islands" and appear as green forested islands in the middle of a sea of desert grasslands and shrub-lands (\cite{moore2012}). Water scarcity is well-documented in this area, and human presence has relied on various forms of water management (\cite{le_tourneau_socio-ecological_2024}), resulting in unique anthropogenic oases scattered throughout the AOI. 

Although the lower concentration levels fall below the EPA's standard of 70 ppb, these same areas experience a heightened risk of DAMO$_3$ during spring and summer as vegetation begins to thrive. This is likely due to CO$_2$ from healthy vegetation found near Pinal farmlands. Additionally, the boundary of these sources primarily comprises areas devoid of vegetation due to the natural desert-like environment present throughout the AOI.Each ArcGIS Pro-based map represents the first day of the month for the corresponding year, as indicated by the census data.
\begin{figure}[h!]
    \centering
    \includegraphics[width=0.8\linewidth]{ThesisFilePkg//figs/maps/dem_o3_2019.png}
    \caption{Averaged 2019 Census Tract Statistics and DAMO$_3$ for January 1st, 2019}
    \label{fig:dem_o3_2019}
\end{figure} In these census-based DAMO$_3$ maps, the transition to a red representation results in yellow, green, and blue shifting to shades of gold-orange, gold, and purple, respectively. Overall, ten maps were created using this method, in fine detail, to demonstrate the potential of this work. 
%%%%%%%%%%%%%%%%%%%%%%%%%%%%%%%%%%%%%%%%%%%%%%%%%%%%%%%%%%%%%%%%
\subsection{Maricopa} \label{discussion:Maricopa} Distributions of DAMO$_3$ for Maricopa County were the most spread out with values ranging from 8.78 ppm to 107.09 ppm. This county houses Phoenix, the most populous city in Arizona as per census. The estimated concentrations ranged from 10.11 ± 3.4 ppb to 91.67 ± 3.4 ppb. The distribution of highly populous areas with smaller amounts of occupied homes can be seen closer to Peoria and Glendale west of Phoenix. The coloring scheme shows many green shaded census tracks near the city center, typically surrounded by counties which have little variations in their mean and median income. Noting the DAMO$_3$ maps, the urban setting in Maricopa exhibits the most diversity in concentration distributions. 
\begin{figure}[h!]
    \centering
    \includegraphics[width=0.8\linewidth]{ThesisFilePkg/figs/Maricopa_ot.png}
    \caption{Average DAMO$_3$ across Pinal County}
    \label{fig:theoryotMaricopa}
\end{figure} When overlaid with corresponding demographic information, many census tracts with high populations, a large number of occupied households, and low deviations in their overall mean versus median income exhibit elevated concentrations. Additionally, during low concentration months such as January, many trends exhibit patterns similar to weather phenomena that tend to encircle certain communities in the area. These same patterns can be seen during the high concentration months like October as well. The effects of urban heat island phenomena and high-quality vegetative areas demonstrate the trends discussed in Chapter \ref{litreview:sofat}, where O$_3$ tends to settle in dry, under-maintained areas, facilitating reactions due to the surrounding environment.

Areas to the Northwest of Phoenix had a stronger weight towards population per capita with a lower weight in high occupancy housing population per capita. This suggests that people living in these areas are either more densely populated in occupied houses or potentially homeless. To the east and southeast, closer to Tempe and Paradise valley, there is less weight towards population per capita, meaning there are fewer people per occupied household. \begin{figure}[h!]
    \centering
    \includegraphics[width=0.35\linewidth]{ThesisFilePkg/figs/maps/ozone_2022-06-24.png}
    \includegraphics[width=0.35\linewidth]{ThesisFilePkg/figs/maps/ozone_2022-06-25.png}
    \includegraphics[width=0.35\linewidth]{ThesisFilePkg/figs/maps/ozone_2022-06-26.png}
    \includegraphics[width=0.35\linewidth]{ThesisFilePkg/figs/maps/ozone_2022-06-27.png}
    \caption{DAMO$_3$ in PHOTUC from June 24th to June 27th, 2022}
    \label{fig:ozone_2022-06-24to27}
\end{figure} Many of the areas that have an equal proportion of population per capita and occupied households live in larger counties denoting potential suburbs of Phoenix. Many of these areas have scattered communities that exist as common stopping points on interstate highways and freeways. Regions closer to the mountains, such as Scottsdale, Paradise Valley, Rio Verde, and Foothills, typically experience higher concentrations of DAMO$_3$, likely attributed to abundant sources of albedo and fluctuations in vegetation cycles. Low concentration seasons such as Winter and Spring show heavier levels of O$_3$ likely due to natural cycles and high elevations in these areas. During high-concentration seasons, these levels are slightly lower, while certain areas of Phoenix experience exacerbated conditions. Glendale, Youngtown, and Sun City West, which tend to see a brunt of high concentrations during the summer post-production at high elevations.
%%%%%%%%%%%%%%%%%%%%%%%%%%%%%%%%%%%%%%%%%%%%%%%%%%%%%%%%%%%%%%%%
\subsection{Pinal} \label{discussion:Pinal} Pinal is situated between Pima and Maricopa, housing the Santa Catalina Mountain range. Many of these communities have the same distribution of population per capita and occupied housing. \begin{figure}[h!]
    \centering
    \includegraphics[width=0.8\linewidth]{ThesisFilePkg/figs/Pinal_ot.png}
    \caption{Average DAMO$_3$ across Pinal County}
    \label{fig:theoryotPinal}
\end{figure} In terms of income, they have seen linear growth and consistent development. The cities and suburbs of Pima see concentrations of ozone from numerous sources during high production seasons. During low concentration seasons, DAMO$_3$ can be seen settling closer towards the Southeast of Maricopa. North of Tucson and Southeast of Phoenix lies Casa Grande, an old mining town located between Phoenix and Tucson. 

Interestingly, Maricopa City is within Pinal. Numerous areas are dedicated to golf courses and farmland, which reflect anthropogenic activity through the types of flora they introduce to the surrounding environment. \begin{figure}[h!]
    \centering
    \includegraphics[width=0.35\linewidth]{ThesisFilePkg/figs/maps/ozone_2020-10-15.png}
    \includegraphics[width=0.35\linewidth]{ThesisFilePkg/figs/maps/ozone_2020-10-16.png}
    \includegraphics[width=0.35\linewidth]{ThesisFilePkg/figs/maps/ozone_2020-10-17.png}
    \includegraphics[width=0.35\linewidth]{ThesisFilePkg/figs/maps/ozone_2020-10-18.png}
    \caption{DAMO$_3$ in PHOTUC from October 1st to October 4th, 2020}
    \label{fig:ozone_2020-10-15to18}
\end{figure} This county is located between the northeastern part of Pima County and the southeastern part of Maricopa County. According to the associated demographic map \ref{appmap:demo3_2020}, many of the purple shaded areas exist East of both major cities, likely due to the large amount of retirement homes. The region also demonstrates significant growth to the southeast of Phoenix and north of Tucson, particularly near Mammoth and Saddlebrook. With the smallest population count, Pinal County exhibits the least overall variation in DAMO $_3$, likely due to extensive anthropogenic vegetation and investments in retirement communities.
%%%%%%%%%%%%%%%%%%%%%%%%%%%%%%%%%%%%%%%%%%%%%%%%%%%%%%%%%%%%%%%%
\subsection{Pima} \label{discussion:Pima} Pima county is one of the original four counties in Arizona and is the oldest of the three in this thesis. Located in Southeastern Arizona, Pima County is characterized by variable climatic conditions, high endemic biodiversity, and rapid urbanization. The region’s history and development are described by a wealth of ecological and social studies, going back to the early days of Native and European settlement \cite{sheridan2012}. While known mainly for the University of Arizona, Distributions of DAMO$_3$ for Pima county were the most spread out with values ranging from 8.78 ppm to 107.09 ppm. The estimated concentrations ranged from 10.11 ± 3.4 ppb to 91.67 ± 3.4 ppb, based on gradient-boosted fields. 
\begin{figure}[h!]
    \centering
    \includegraphics[width=0.8\linewidth]{ThesisFilePkg/figs/Pima_ot.png}
    \caption{Average DAMO$_3$ across Pima County}
    \label{fig:theoryotPima}
\end{figure} When overlaid with corresponding demographic information, some census tracts with high populations, a low number of occupied households, and minimal deviations between mean and median values, as well as average median income, exhibit higher concentrations during the summer months. Additionally, during low concentration months such as January, many trends align with weather patterns that tend to encircle certain communities, particularly near Tucson at the base of Mt. Lemmon. Similar patterns are also observable during high concentration months, such as October. The impacts of both urban heat island effects and high-quality vegetative areas reflect the trends discussed in Chapter \ref{chap:litreview} where O$_3$ tends to settle in dry, under-maintained areas, preventing interaction with vegetation.

Areas to the southwest of Phoenix emphasized population density more significantly, as indicated by numerous gold-shaded census tracts, while placing less emphasis on high-occupancy housing population per capita. In contrast, to the east and southeast, closer to Tempe and Paradise Valley, there is a reduced emphasis on population density. This area exhibits fewer people per occupied household, coupled with higher average incomes. Regions where there is an equal proportion of population density and occupied households are predominantly located in larger counties, suggesting they may be potential suburbs of Pima. Additionally, many of these areas feature scattered communities that serve as common stopping points along interstate highways and freeways. 
\begin{figure}[h!]
    \centering
    \includegraphics[width=0.35\linewidth]{ThesisFilePkg/figs/maps/ozone_2019-01-10.png}
    \includegraphics[width=0.35\linewidth]{ThesisFilePkg/figs/maps/ozone_2019-01-11.png}
    \includegraphics[width=0.35\linewidth]{ThesisFilePkg/figs/maps/ozone_2019-01-12.png}
    \includegraphics[width=0.35\linewidth]{ThesisFilePkg/figs/maps/ozone_2019-01-13.png}
    \caption{DAMO$_3$ in PHOTUC from January 10th to January 13th, 2019}
    \label{fig:ozone_2019-01-10to13}
\end{figure} Areas closer to the mountains, such as Scottsdale, Paradise Valley, Rio Verde, and the Foothills, tend to experience higher concentrations of DAMO$_3$, likely due to increased levels of albedo and vegetation cycles. Low concentration seasons such as Winter and Spring show heavier levels of O$_3$ likely due to natural cycles and high elevations in these areas. During high concentration seasons, these are slightly lower while exacerbating certain areas of Phoenix. Glendale, Youngtown, and Sun City bear the brunt of high concentrations during the summer, particularly at elevated positions. The University of Arizona is within Tucson, along with a heavy Raytheon presence, entities which contribute to healthier vegetation and better care of either natural or invasive flora. Many of the higher concentrations of ozone can be seen near Superior and Casa Grande in \ref{fig:ozone_2019-01-10to13} due to the lack of driving forces in the winter.
%%%%%%%%%%%%%%%%%%%%%%%%%%%%%%%%%%%%%%%%%%%%%%%%%%%%%%%%%%%%%%%%%%%%%%%%%%%%%%%%%%%%%%%%%%%%%%%%%%%%%%%%%%%%%%%%%%%%%%%%%%%%%%%%
\section{Final Trends Across PHOTUC} \label{discussion:FinConcTrends} Areas with little green space near both Phoenix and Tucson would experience the highest concentrations during summer. Average levels, which exceed 70 ppb according to both in-situ and modeled measurements, predominantly occur in densely populated areas surrounding Phoenix and Tucson during this season as well. Many census tracts exhibit significant variations in DAMO$_3$ within and just outside of these urban classified boundaries. By effectively addressing spatial and temporal variability, RK can support targeted interventions to reduce DAMO$_3$ exposure in dense areas such as PHOTUC. 

Policies informed by high-resolution DAMO$_3$ models can guide infrastructure development and resource allocation for air quality improvement initiatives, particularly in vulnerable urban populations. The capacity of RK to capture fine-scale DAMO$_3$ variations has the potential to advance urban planning by better identifying trends across census tracts, where localized emission control measures could significantly impact public health. The successful integration of RK into XGRB, along with the accompanying suite of complex modeling ensembles, is further elucidated through the supplementary imagery provided in the Appendix. This study’s findings underscore the necessity of high-resolution models for achieving accurate exposure assessments that inform public health strategies and urban expansion plans. 

The ecosystems in the AOI interact with DAMO$_3$ throughout each monthly depiction of DAMO$_3$. The final daily images provided on Map \ref{fig:ozone_2021-07-28to31} show elevation spikes tend to trap reactions within lower elevations after generation near urban areas. This is offset by areas with healthy vegetation. During Fall, significant changes in DAMO$_3$ concentrations can be observed in agricultural regions where vegetation is preparing for winter, particularly east of Casa Grande and south of Blackwater. In addition, sky islands such as Mount Lemmon and other spikes in elevation, tend to have similar cycles due to their ability to contain difference biomes \parencite{moore2012}. Such changes can be seen in \ref{fig:ozone_2021-07-28to31} where most of the DAMO$_3$ concentration settles outside of Tucson's boundary.
\begin{figure}[h!]
    \centering
    \includegraphics[width=0.35\linewidth]{ThesisFilePkg/figs/maps/ozone_2021-07-28.png}
    \includegraphics[width=0.35\linewidth]{ThesisFilePkg/figs/maps/ozone_2021-07-29.png}
    \includegraphics[width=0.35\linewidth]{ThesisFilePkg/figs/maps/ozone_2021-07-30.png}
    \includegraphics[width=0.35\linewidth]{ThesisFilePkg/figs/maps/ozone_2021-07-31.png}
    \caption{DAMO$_3$ in PHOTUC from July 4th to July 7th, 2021}
    \label{fig:ozone_2021-07-28to31}
\end{figure} Areas near the basins of the Rincon and Santa Catalina mountain ranges would also see prolonged after-effects after ozone production, likely due to nearby developing areas. During winter seasons, low to mid-estimates of ozone can be seen settling in the suburbs of both Phoenix and Tucson. 
%%%%%%%%%%%%%%%%%%%%%%%%%%%%%%%%%%%%%%%%%%%%%%%%%%%%%%%%%%%%%%%%%%%%%%%%%%%%%%%%%%%%%%%%%%%%%%%%%%%%%%%%%%%%%%%%%%%%%%%%%%%%%%%%
\subsection{Potential Impacts of DAMO$_3$ Modeling} \label{discussion:histtrends} The historical feasibility of high-resolution DAMO$_3$ modeling, particularly from 1980 to 2000, relies heavily on the RK method’s ability to accurately capture residual spatial dependencies. With RK, historical data from sources like TOMS can be spatially refined to produce monthly average DAMO$_3$ concentrations, filling a critical gap in historical DAMO$_3$ exposure records. This capability is invaluable for public health research, as it facilitates a retrospective analysis of exposure patterns over time, establishing connections between long-term DAMO$_3$ trends and health outcomes across various populations. The RK method’s refinement of historical datasets supports an unprecedented level of detail in exposure mapping, fostering a deeper understanding of how DAMO$_3$ exposure has evolved in response to environmental and socioeconomic changes. From 2000 onward, with the availability of advanced satellite data from instruments such as S5P, the RK method facilitates daily average DAMO$_3$ models through granular spatial predictions, making high-resolution daily models more feasible.

One of the most promising aspects of the RK method lies in its potential to improve existing statistical models by estimating predictions prior to 2005. By predicting geospatial error over a broader temporal range, RK can enhance the predictive accuracy of ML/AI corrections, which are used to understand trends in collected imagery spanning several decades. When combined with residual kriging, statistical models can deliver predictions that reflect both large-scale trends and localized variations. This integration fosters a new generation of high-resolution spatial datasets that provide valuable insights across numerous fields, enabling more precise estimations of the impacts related to different DAMO$_3$ exposure levels and beyond. The RK method offers a systematic approach to account for the inherent variability in DAMO$_3$ concentrations, which traditional statistical models often struggle to address. This adjustment is critical for capturing DAMO$_3$’s complex behavior across diverse landscapes, which is essential for accurate long-term health impact assessments.